\documentclass[12pt,A4, onecolumn]{exam}
\usepackage{amsmath}
\usepackage{amssymb}
\usepackage[lmargin=71pt, tmargin=1.2in]{geometry}  %For centering solution box
\lhead{TEK4090 -- Introduction to Modern Control\\}
\rhead{Assignment 1\\}
% \chead{\hline} % Un-comment to draw line below header
\thispagestyle{empty}   %For removing header/footer from page 1

\usepackage{graphicx}
\graphicspath{{../../figs/}}

\def\showsoln{1}

\begin{document}

\begingroup  
    \centering
    \LARGE TEK4090 -- Introduction to Modern Control\\
    \LARGE Assignment 1\\[0.5em]
    \large Due: September 18 2025, 13:15\\
\endgroup
\rule{\textwidth}{0.4pt}
%\pointsdroppedatright   %Self-explanatory
\printanswers


\renewcommand{\solutiontitle}{\noindent\textbf{Ans:}\enspace}   %Replace "Ans:" with starting keyword in solution box

\begin{questions}

  \question[10] Suppose that $v_1,\dots,v_m$ is linearly independent in $V$, and that $w\in V$. Show that
    \begin{align}
      v_1,\dots,v_m, w \text{ is linearly independent } \iff w\notin \mathrm{span}(v_1,\dots,v_m).
    \end{align}


  \question[10] Let $A,B\in \mathbb{R}^{n\times n}$, and define the \emph{trace} of $A$ as $\mathrm{tr}(A)=\sum_{i=1}^n A_{ii}$. Show that 
    \begin{align}
      \dfrac{d \mathrm{tr}(AB)}{dA} = B^T.
    \end{align}
    Hint: Find $\frac{d\mathrm{tr}(AB)}{dA_{ij}}$.

  \question[15] Suppose that $T:V\mapsto V$ is a linear map, and suppose that $S:V\mapsto V$ is an invertible linear map.
    \begin{parts}
      \part Prove that $T$ and $S^{-1}TS$ have the same eigenvalues
      \part Describe the relationship between the eigenvectors of $T$ and the eigenvectors of $S^{-1}TS$.
    \end{parts}

\question[30] Consider the feedback block diagram:
  \begin{center}
    \includegraphics[width=0.75\linewidth]{./feedback}
  \end{center}
\begin{parts}
  \part Find the transfer function from $d$ to $y$.
  \part Suppose the plant has transfer function
\begin{align}
  P(s) = \dfrac{(s+1)}{s^2-4s+13}
\end{align} 
What are the poles of $P(s)$? What can you infer about the stability of the (open-loop) plant $P(s)$?
  \part Suppose $C(s)$ is given by 
    \begin{align}
      C(s) = \dfrac{1}{s}
    \end{align}
    Does $C(s)$ stabilize the feedback system?  
  \part \label{1} Now, suppose the plant and controller are given by
    \begin{align}
      P(s) = \dfrac{10}{(s+1)},~C(s) =\dfrac{K}{s}.
    \end{align}
  Find the values of $K$ for which the feedback loop is stable.
  \part Fix a value of $K$ such that the system in part~(\ref{1}) is stable, and suppose the system is subject to a disturbance $D(s) = 100$. What will the final value of $y$ be?
\end{parts}

    

    \question[10] Consider the standard feedback system with
\begin{align}
C(s) = K,~
P (s) =\dfrac{
s^2 + 1}{
(s^2 + 2)(s - 2)}
\end{align}
\begin{parts}
  
\part Sketch the Nyquist plot of $P (s)$. Include correct asymptotes to infinity and to the origin.
\part Indicate the number of encirclements (and their orientations) of all possible critical points.
\part For what range of $K$ is the feedback system stable?
\end{parts}


  \begin{center}
    \includegraphics[width=0.6\linewidth]{./bode-hwk-1}\\
    \textbf{Fig.~1:} Bode Plot for Question~(\ref{q:bode})
  \end{center}

\begin{question}[15]
  \label{q:bode} Consider the standard feedback system with $C(s) = K$. 
  The Bode plot of $P (s)$ is given in Figure 1 (magnitude in absolute units, not dB; phase in degrees). 
  The phase starts at $-180^\circ$ and ends at $-270^\circ$. 
  You are also given that $P(s)$ has exactly one pole in the right half-plane. 
  For what range of gains $K$ is the feedback system stable?
\end{question}

\question[10] Consider the feedback system:
\begin{center}
  \includegraphics[width=0.5\linewidth]{../figs/k-system}
\end{center}
with
\begin{align}
  P(s) =
    \dfrac{s+2}{s^2+2},~C(s) = \dfrac{1}{s}.
\end{align}

Sketch the Nyquist plot of $PC$. How many encirclements are required of the critical point for feedback stability? Determine the range of real gains $K$ for stability of the feedback system.


\question[10]

Repeat the previous question with,
\begin{align}
  P(s) = \dfrac{s^2+1}{(s+1)(s^2+s+1)},~C(s) = 1.
\end{align}

\end{questions}


\end{document}